%%
% This is an Overleaf template for scientific posters
% using the TUM Corporate Desing https://www.tum.de/cd
%
% For further details on how to use the template, take a look at our
% GitLab repository and browse through our test documents
% https://gitlab.lrz.de/latex4ei/tum-templates.
%
% The tumposter class is based on the KOMA-Script class scrartcl and the
% poster framework of the package tcolorbox.
% If you need further customization please consult the KOMA-Script guide
% https://ctan.org/pkg/koma-script
% and the tcolorbox user guide (especially the poster section)
% https://ctan.org/pkg/tcolorbox.
% Additional class options are passed down to the base class.
%
% If you encounter any bugs or undesired behaviour, please raise an issue
% in our GitLab repository
% https://gitlab.lrz.de/latex4ei/tum-templates/issues
% and provide a description and minimal working example of your problem.
%%


\documentclass[
  a0paper,    % poster size (a0paper, a1paper, a2paper, a3paper, a4paper)
  portrait,   % poster orientation (portrait, landscape)
  english,    % define the document language (english, german)
]{tumposter}


% load additional packages
\usepackage{tummath}
\usepackage{lipsum}


% poster metadata
\title{Poster title}
\subtitle{Subtitle of the poster}

\author[affil=1]{First Author}
\author[affil=2]{Second Author}
\author[affil=1]{Last Author}

\affil[mark=1]{\theSchoolName, \theUniversityName}
\affil[mark=2]{Another affiliation}


% macro to configure the style of the poster
\TUMpostersetup{
  columns = 4,            % number of columns of the poster
  rows = 6,               % number of rows of the poster
  background = white,     % background color (white, blue, shaded)
  boxes = TUMBoxBlank,    % default box style (TUMBox, TUMBoxVariant, TUMBoxInverse, TUMBoxBlank)
  headline = threeliner,  % choose headline (tumlogo, oneliner, twoliner, threeliner)
  % printaffil = false,     % disable affiliation in title area (in case you print it in a footer)
  % fontsize = 25pt,        % choose larger base font size if necessary
  % showframe = true,       % print poster row and column grid (for draft versions)
}


\begin{document}

\maketitle

\begin{tcbposter}

  % footer with some contact information
  \posterbox[fontupper=\small, height=5cm,
    enlarge top initially by = .5\TUMposterboxsep
  ]{
    name=footer, column=1, span=4, above=bottom
  }{
    \textbf{\theUniversityName}\\
    This footer box could be used to place contact information.\\
    Setting a fixed height and insert additional spacing\\
    to separate the footer from the main content is possible.
  }

  % figure on the top right
  \posterbox{name=fig1, column=4, below=top}{
    \includegraphics[width=\columnwidth, height=30cm]{example-image-a}
    \captionof{figure}{This is a figure caption. The figure contains a
      placeholder to test the layout.}
  }

  % figure on the bottom left
  \posterbox{name=fig2, column=1, above=footer}{
    \includegraphics[width=\columnwidth, height=20cm]{example-image-b}
    \captionof{figure}{This is another figure placeholder.}
  }

  % main text box is a sequence over the rest of the page
  \posterbox{name=text, column=1,
    sequence = 1 between top and fig2 then
    2 between top and footer then
    3 between top and footer then
    4 between fig1 and footer
  }{
    \lipsum[1-15]
  }

\end{tcbposter}

\end{document}
